% Paper A: proofs for Section 3 and Section 4.
% Included by paper.tex in the appendix (% Paper A: proofs for Section 3 and Section 4.
% Included by paper.tex in the appendix (% Paper A: proofs for Section 3 and Section 4.
% Included by paper.tex in the appendix (% Paper A: proofs for Section 3 and Section 4.
% Included by paper.tex in the appendix (\input{paper_a_proofs}).
% Theorem 1's proof stays in the main text: the later proofs and the cut
% argument of Section 4 all cite "the argument of Theorem 1".
% Check status follows the \status tag on each statement in the main text.

\subsection{Proofs for Section~\ref{sec:guarantee}}

\begin{proof}[Proof of Corollary~\ref{cor:finite}]
Let $s_1,\dots,s_{K'}$, $K' \le K$, be the distinct values with
$\|s_k\|_\Sigma \neq 0$, and $u_k = \Sigma^{1/2}s_k / \|s_k\|_\Sigma$, so that
$W = \max_k \, -u_k^\top\xi$. Each term is sub-Gaussian with variance proxy
$\|u_k\|_2^2 = 1$ \citep[Proposition~15.2]{bertsimas2022robust}, so
$\mathbb{P}(-u_k^\top\xi > R) \le \exp(-R^2/2)$
\citep[Proposition~15.1]{bertsimas2022robust}. A union bound gives
$\mathbb{P}(W > R) \le K\exp(-R^2/2) = \epsilon$.
\end{proof}

\begin{proof}[Proof of Corollary~\ref{cor:subspace}]
Let $\Pi$ be the orthogonal projection onto $\Sigma^{1/2}S$, a subspace of
dimension $d' \le d$. Since $\Sigma^{1/2}s(x) \in \Sigma^{1/2}S$,
$-s(x)^\top\varepsilon = -(\Sigma^{1/2}s(x))^\top \Pi\xi
\le \|s(x)\|_\Sigma\,\|\Pi\xi\|_2$, so $W \le \|\Pi\xi\|_2$. Moreover
$\|\Pi\xi\|_2^2 \sim \chi^2_{d'}$, which is stochastically no larger than
$\chi^2_d$, hence
$\mathbb{P}(W > R) \le \mathbb{P}(\|\Pi\xi\|_2^2 > \chi^2_{d,1-\epsilon})
\le \epsilon$.
\end{proof}

\begin{proof}[Proof of Corollary~\ref{cor:multi}]
By the argument of Theorem~\ref{thm:selection}, the event that $x^\star$
violates some constraint is contained in $\bigcup_{j=1}^m \{W_j > R_j\}$.
Each event has probability at most $\epsilon/m$; apply a union bound.
\end{proof}

\begin{proof}[Proof of Proposition~\ref{prop:simulated}]
The condition on $N$ gives $k \le N$. Let $W_i = w(\xi^{(i)})$ and let
$W_{(k)}$ be the $k$-th smallest of $W_1,\dots,W_N$. The argument of
Theorem~\ref{thm:selection} holds pointwise for every value of the radius, so
$f(x^\star;\theta^\star) > b$ implies $W > \hat R$. Since
$\overline W_i \ge W_i$ for every $i$, $\hat R \ge W_{(k)}$, and it suffices to
show $\mathbb{P}(W > W_{(k)}) \le \epsilon$.

The variables $V_0 = W$ and $V_i = W_i$, $i = 1,\dots,N$, are independent and
identically distributed, hence exchangeable, and $W > W_{(k)}$ holds exactly
when $V_0$ is strictly larger than at least $k$ of the other $N$. Order
$V_0,\dots,V_N$ increasingly, breaking ties arbitrarily. The variable in
position $j$ is strictly larger than at most $j-1$ of the others, so at most
$N+1-k$ of the $N+1$ variables are strictly larger than at least $k$ of the
others. By exchangeability each variable has the same probability of being
one of them, so
$\mathbb{P}(W > W_{(k)}) \le (N+1-k)/(N+1) \le \epsilon$, the last step
because $k \ge (N+1)(1-\epsilon)$.
\end{proof}

\subsection{Proof for Section~\ref{sec:solving}}

\begin{proof}[Proof of Proposition~\ref{prop:cuts}]
By the Cauchy--Schwarz inequality in the inner product
$\langle u, v\rangle = u^\top\Sigma v$, the left side of \eqref{eq:cut} is at
most $f(x;\hat\theta) + R\|s(x)\|_\Sigma$, so every cut is implied by the
robust counterpart and every master problem is a relaxation of the robust
problem. If $\|s^k\|_\Sigma = 0$, the counterpart at $x^k$ is the nominal
constraint, which the master problem enforces, so cuts are added only at
iterates with $\|s^k\|_\Sigma \neq 0$. Such an iterate has a value of $s$ not
seen at any earlier iterate: had $s^k = s^j$ for some $j < k$, cut $j$ would
read $f(x^k;\hat\theta) + R\|s^k\|_\Sigma \le b$ at $x^k$, which is the robust
counterpart. So each cut is added at a new value of $s$, and there are at
most $K$. The point returned satisfies the robust counterpart and is optimal
for a relaxation of the robust problem, hence optimal for it.
\end{proof}
).
% Theorem 1's proof stays in the main text: the later proofs and the cut
% argument of Section 4 all cite "the argument of Theorem 1".
% Check status follows the \status tag on each statement in the main text.

\subsection{Proofs for Section~\ref{sec:guarantee}}

\begin{proof}[Proof of Corollary~\ref{cor:finite}]
Let $s_1,\dots,s_{K'}$, $K' \le K$, be the distinct values with
$\|s_k\|_\Sigma \neq 0$, and $u_k = \Sigma^{1/2}s_k / \|s_k\|_\Sigma$, so that
$W = \max_k \, -u_k^\top\xi$. Each term is sub-Gaussian with variance proxy
$\|u_k\|_2^2 = 1$ \citep[Proposition~15.2]{bertsimas2022robust}, so
$\mathbb{P}(-u_k^\top\xi > R) \le \exp(-R^2/2)$
\citep[Proposition~15.1]{bertsimas2022robust}. A union bound gives
$\mathbb{P}(W > R) \le K\exp(-R^2/2) = \epsilon$.
\end{proof}

\begin{proof}[Proof of Corollary~\ref{cor:subspace}]
Let $\Pi$ be the orthogonal projection onto $\Sigma^{1/2}S$, a subspace of
dimension $d' \le d$. Since $\Sigma^{1/2}s(x) \in \Sigma^{1/2}S$,
$-s(x)^\top\varepsilon = -(\Sigma^{1/2}s(x))^\top \Pi\xi
\le \|s(x)\|_\Sigma\,\|\Pi\xi\|_2$, so $W \le \|\Pi\xi\|_2$. Moreover
$\|\Pi\xi\|_2^2 \sim \chi^2_{d'}$, which is stochastically no larger than
$\chi^2_d$, hence
$\mathbb{P}(W > R) \le \mathbb{P}(\|\Pi\xi\|_2^2 > \chi^2_{d,1-\epsilon})
\le \epsilon$.
\end{proof}

\begin{proof}[Proof of Corollary~\ref{cor:multi}]
By the argument of Theorem~\ref{thm:selection}, the event that $x^\star$
violates some constraint is contained in $\bigcup_{j=1}^m \{W_j > R_j\}$.
Each event has probability at most $\epsilon/m$; apply a union bound.
\end{proof}

\begin{proof}[Proof of Proposition~\ref{prop:simulated}]
The condition on $N$ gives $k \le N$. Let $W_i = w(\xi^{(i)})$ and let
$W_{(k)}$ be the $k$-th smallest of $W_1,\dots,W_N$. The argument of
Theorem~\ref{thm:selection} holds pointwise for every value of the radius, so
$f(x^\star;\theta^\star) > b$ implies $W > \hat R$. Since
$\overline W_i \ge W_i$ for every $i$, $\hat R \ge W_{(k)}$, and it suffices to
show $\mathbb{P}(W > W_{(k)}) \le \epsilon$.

The variables $V_0 = W$ and $V_i = W_i$, $i = 1,\dots,N$, are independent and
identically distributed, hence exchangeable, and $W > W_{(k)}$ holds exactly
when $V_0$ is strictly larger than at least $k$ of the other $N$. Order
$V_0,\dots,V_N$ increasingly, breaking ties arbitrarily. The variable in
position $j$ is strictly larger than at most $j-1$ of the others, so at most
$N+1-k$ of the $N+1$ variables are strictly larger than at least $k$ of the
others. By exchangeability each variable has the same probability of being
one of them, so
$\mathbb{P}(W > W_{(k)}) \le (N+1-k)/(N+1) \le \epsilon$, the last step
because $k \ge (N+1)(1-\epsilon)$.
\end{proof}

\subsection{Proof for Section~\ref{sec:solving}}

\begin{proof}[Proof of Proposition~\ref{prop:cuts}]
By the Cauchy--Schwarz inequality in the inner product
$\langle u, v\rangle = u^\top\Sigma v$, the left side of \eqref{eq:cut} is at
most $f(x;\hat\theta) + R\|s(x)\|_\Sigma$, so every cut is implied by the
robust counterpart and every master problem is a relaxation of the robust
problem. If $\|s^k\|_\Sigma = 0$, the counterpart at $x^k$ is the nominal
constraint, which the master problem enforces, so cuts are added only at
iterates with $\|s^k\|_\Sigma \neq 0$. Such an iterate has a value of $s$ not
seen at any earlier iterate: had $s^k = s^j$ for some $j < k$, cut $j$ would
read $f(x^k;\hat\theta) + R\|s^k\|_\Sigma \le b$ at $x^k$, which is the robust
counterpart. So each cut is added at a new value of $s$, and there are at
most $K$. The point returned satisfies the robust counterpart and is optimal
for a relaxation of the robust problem, hence optimal for it.
\end{proof}
).
% Theorem 1's proof stays in the main text: the later proofs and the cut
% argument of Section 4 all cite "the argument of Theorem 1".
% Check status follows the \status tag on each statement in the main text.

\subsection{Proofs for Section~\ref{sec:guarantee}}

\begin{proof}[Proof of Corollary~\ref{cor:finite}]
Let $s_1,\dots,s_{K'}$, $K' \le K$, be the distinct values with
$\|s_k\|_\Sigma \neq 0$, and $u_k = \Sigma^{1/2}s_k / \|s_k\|_\Sigma$, so that
$W = \max_k \, -u_k^\top\xi$. Each term is sub-Gaussian with variance proxy
$\|u_k\|_2^2 = 1$ \citep[Proposition~15.2]{bertsimas2022robust}, so
$\mathbb{P}(-u_k^\top\xi > R) \le \exp(-R^2/2)$
\citep[Proposition~15.1]{bertsimas2022robust}. A union bound gives
$\mathbb{P}(W > R) \le K\exp(-R^2/2) = \epsilon$.
\end{proof}

\begin{proof}[Proof of Corollary~\ref{cor:subspace}]
Let $\Pi$ be the orthogonal projection onto $\Sigma^{1/2}S$, a subspace of
dimension $d' \le d$. Since $\Sigma^{1/2}s(x) \in \Sigma^{1/2}S$,
$-s(x)^\top\varepsilon = -(\Sigma^{1/2}s(x))^\top \Pi\xi
\le \|s(x)\|_\Sigma\,\|\Pi\xi\|_2$, so $W \le \|\Pi\xi\|_2$. Moreover
$\|\Pi\xi\|_2^2 \sim \chi^2_{d'}$, which is stochastically no larger than
$\chi^2_d$, hence
$\mathbb{P}(W > R) \le \mathbb{P}(\|\Pi\xi\|_2^2 > \chi^2_{d,1-\epsilon})
\le \epsilon$.
\end{proof}

\begin{proof}[Proof of Corollary~\ref{cor:multi}]
By the argument of Theorem~\ref{thm:selection}, the event that $x^\star$
violates some constraint is contained in $\bigcup_{j=1}^m \{W_j > R_j\}$.
Each event has probability at most $\epsilon/m$; apply a union bound.
\end{proof}

\begin{proof}[Proof of Proposition~\ref{prop:simulated}]
The condition on $N$ gives $k \le N$. Let $W_i = w(\xi^{(i)})$ and let
$W_{(k)}$ be the $k$-th smallest of $W_1,\dots,W_N$. The argument of
Theorem~\ref{thm:selection} holds pointwise for every value of the radius, so
$f(x^\star;\theta^\star) > b$ implies $W > \hat R$. Since
$\overline W_i \ge W_i$ for every $i$, $\hat R \ge W_{(k)}$, and it suffices to
show $\mathbb{P}(W > W_{(k)}) \le \epsilon$.

The variables $V_0 = W$ and $V_i = W_i$, $i = 1,\dots,N$, are independent and
identically distributed, hence exchangeable, and $W > W_{(k)}$ holds exactly
when $V_0$ is strictly larger than at least $k$ of the other $N$. Order
$V_0,\dots,V_N$ increasingly, breaking ties arbitrarily. The variable in
position $j$ is strictly larger than at most $j-1$ of the others, so at most
$N+1-k$ of the $N+1$ variables are strictly larger than at least $k$ of the
others. By exchangeability each variable has the same probability of being
one of them, so
$\mathbb{P}(W > W_{(k)}) \le (N+1-k)/(N+1) \le \epsilon$, the last step
because $k \ge (N+1)(1-\epsilon)$.
\end{proof}

\subsection{Proof for Section~\ref{sec:solving}}

\begin{proof}[Proof of Proposition~\ref{prop:cuts}]
By the Cauchy--Schwarz inequality in the inner product
$\langle u, v\rangle = u^\top\Sigma v$, the left side of \eqref{eq:cut} is at
most $f(x;\hat\theta) + R\|s(x)\|_\Sigma$, so every cut is implied by the
robust counterpart and every master problem is a relaxation of the robust
problem. If $\|s^k\|_\Sigma = 0$, the counterpart at $x^k$ is the nominal
constraint, which the master problem enforces, so cuts are added only at
iterates with $\|s^k\|_\Sigma \neq 0$. Such an iterate has a value of $s$ not
seen at any earlier iterate: had $s^k = s^j$ for some $j < k$, cut $j$ would
read $f(x^k;\hat\theta) + R\|s^k\|_\Sigma \le b$ at $x^k$, which is the robust
counterpart. So each cut is added at a new value of $s$, and there are at
most $K$. The point returned satisfies the robust counterpart and is optimal
for a relaxation of the robust problem, hence optimal for it.
\end{proof}
).
% Theorem 1's proof stays in the main text: the later proofs and the cut
% argument of Section 4 all cite "the argument of Theorem 1".
% Check status follows the \status tag on each statement in the main text.

\subsection{Proofs for Section~\ref{sec:guarantee}}

\begin{proof}[Proof of Corollary~\ref{cor:finite}]
Let $s_1,\dots,s_{K'}$, $K' \le K$, be the distinct values with
$\|s_k\|_\Sigma \neq 0$, and $u_k = \Sigma^{1/2}s_k / \|s_k\|_\Sigma$, so that
$W = \max_k \, -u_k^\top\xi$. Each term is sub-Gaussian with variance proxy
$\|u_k\|_2^2 = 1$ \citep[Proposition~15.2]{bertsimas2022robust}, so
$\mathbb{P}(-u_k^\top\xi > R) \le \exp(-R^2/2)$
\citep[Proposition~15.1]{bertsimas2022robust}. A union bound gives
$\mathbb{P}(W > R) \le K\exp(-R^2/2) = \epsilon$.
\end{proof}

\begin{proof}[Proof of Corollary~\ref{cor:subspace}]
Let $\Pi$ be the orthogonal projection onto $\Sigma^{1/2}S$, a subspace of
dimension $d' \le d$. Since $\Sigma^{1/2}s(x) \in \Sigma^{1/2}S$,
$-s(x)^\top\varepsilon = -(\Sigma^{1/2}s(x))^\top \Pi\xi
\le \|s(x)\|_\Sigma\,\|\Pi\xi\|_2$, so $W \le \|\Pi\xi\|_2$. Moreover
$\|\Pi\xi\|_2^2 \sim \chi^2_{d'}$, which is stochastically no larger than
$\chi^2_d$, hence
$\mathbb{P}(W > R) \le \mathbb{P}(\|\Pi\xi\|_2^2 > \chi^2_{d,1-\epsilon})
\le \epsilon$.
\end{proof}

\begin{proof}[Proof of Corollary~\ref{cor:multi}]
By the argument of Theorem~\ref{thm:selection}, the event that $x^\star$
violates some constraint is contained in $\bigcup_{j=1}^m \{W_j > R_j\}$.
Each event has probability at most $\epsilon/m$; apply a union bound.
\end{proof}

\begin{proof}[Proof of Proposition~\ref{prop:simulated}]
The condition on $N$ gives $k \le N$. Let $W_i = w(\xi^{(i)})$ and let
$W_{(k)}$ be the $k$-th smallest of $W_1,\dots,W_N$. The argument of
Theorem~\ref{thm:selection} holds pointwise for every value of the radius, so
$f(x^\star;\theta^\star) > b$ implies $W > \hat R$. Since
$\overline W_i \ge W_i$ for every $i$, $\hat R \ge W_{(k)}$, and it suffices to
show $\mathbb{P}(W > W_{(k)}) \le \epsilon$.

The variables $V_0 = W$ and $V_i = W_i$, $i = 1,\dots,N$, are independent and
identically distributed, hence exchangeable, and $W > W_{(k)}$ holds exactly
when $V_0$ is strictly larger than at least $k$ of the other $N$. Order
$V_0,\dots,V_N$ increasingly, breaking ties arbitrarily. The variable in
position $j$ is strictly larger than at most $j-1$ of the others, so at most
$N+1-k$ of the $N+1$ variables are strictly larger than at least $k$ of the
others. By exchangeability each variable has the same probability of being
one of them, so
$\mathbb{P}(W > W_{(k)}) \le (N+1-k)/(N+1) \le \epsilon$, the last step
because $k \ge (N+1)(1-\epsilon)$.
\end{proof}

\subsection{Proof for Section~\ref{sec:solving}}

\begin{proof}[Proof of Proposition~\ref{prop:cuts}]
By the Cauchy--Schwarz inequality in the inner product
$\langle u, v\rangle = u^\top\Sigma v$, the left side of \eqref{eq:cut} is at
most $f(x;\hat\theta) + R\|s(x)\|_\Sigma$, so every cut is implied by the
robust counterpart and every master problem is a relaxation of the robust
problem. If $\|s^k\|_\Sigma = 0$, the counterpart at $x^k$ is the nominal
constraint, which the master problem enforces, so cuts are added only at
iterates with $\|s^k\|_\Sigma \neq 0$. Such an iterate has a value of $s$ not
seen at any earlier iterate: had $s^k = s^j$ for some $j < k$, cut $j$ would
read $f(x^k;\hat\theta) + R\|s^k\|_\Sigma \le b$ at $x^k$, which is the robust
counterpart. So each cut is added at a new value of $s$, and there are at
most $K$. The point returned satisfies the robust counterpart and is optimal
for a relaxation of the robust problem, hence optimal for it.
\end{proof}
